\documentclass[../DoAn.tex]{subfiles}
\begin{document}

\begin{center}
    \Large{\textbf{TÓM TẮT NỘI DUNG ĐỒ ÁN}}\\
\end{center}
\vspace{1cm}
% TODO: cập nhật câu cuối khi có kết quả vòng lặp LLM, Optuna và dữ liệu quan sát thực.
Mô phỏng vi mô là công cụ phổ biến để phân tích nút giao có đèn tín hiệu, nhưng các mô hình hành vi mặc định của trình mô phỏng như SUMO được xây dựng cho dòng xe tuân thủ luật và có làn rõ ràng. Giao thông hỗn hợp ô tô -- xe máy tại các đô thị Việt Nam có nhiều hành vi phi chuẩn như vượt đèn đỏ theo đoàn, xuất phát sớm, lấn chiếm theo phương ngang và rẽ trái trái phép bằng cách dừng chờ giữa nút. Việc hiệu chỉnh thủ công một không gian tham số lớn để tái hiện các hành vi này tốn nhiều công sức và khó truy vết. Đồ án đề xuất một quy trình trong đó mô hình ngôn ngữ lớn (LLM) chạy cục bộ sinh và tinh chỉnh lặp một đặc tả luật hành vi tài xế, gồm tham số kiểu phương tiện tĩnh và các luật có điều kiện thực thi qua TraCI; LLM chỉ sinh luật ngoài vòng mô phỏng, còn việc thực thi luật do SUMO và một bộ điều khiển TraCI đảm nhận. Mọi đặc tả được kiểm tra hợp lệ theo lược đồ và miền tham số có ghi nguồn trước khi biên dịch và chạy. Quy trình được áp dụng cho nút Lê Văn Lương -- Hoàng Đạo Thúy (Hà Nội) trong cao điểm sáng, với kịch bản trọng tâm là rẽ trái trái phép. Đến thời điểm hiện tại, đồ án đã xây dựng mạng mô phỏng tái lập được, cấu hình cơ sở, thí nghiệm khả thi về xe chờ giữa nút, bộ kiểm tra hợp lệ, bộ biên dịch và bộ điều khiển với năm luật có điều kiện; hệ thống chỉ số đang được triển khai, còn vòng lặp LLM và phương án đối sánh tối ưu bằng Optuna với cùng ngân sách số lần chạy là phần dự kiến. Do chưa có dữ liệu quan sát từ video, quy trình hiện được kiểm thử trên các ngày dữ liệu đại diện tổng hợp từ tài liệu; kết quả trên dữ liệu này chỉ đánh giá khả năng vận hành của quy trình, không khẳng định mức độ tương đồng với giao thông thực tế. \textbf{[TODO: đóng góp chính và kết quả đạt được]}
\begin{flushright}
Sinh viên thực hiện\\
\begin{tabular}{@{}c@{}}
\textit{(Ký và ghi rõ họ tên)}
\end{tabular}
\end{flushright}

\end{document}
